\documentclass[11pt]{article}

\usepackage[T1]{fontenc}
\usepackage{amsmath,amssymb,amsthm}
\usepackage[margin=1in]{geometry}
\usepackage[hidelinks]{hyperref}
\setlength{\emergencystretch}{3em}

\newtheorem{theorem}{Theorem}
\newtheorem{lemma}[theorem]{Lemma}
\newtheorem{corollary}[theorem]{Corollary}
\newtheorem{proposition}[theorem]{Proposition}
\newtheorem{conjecture}[theorem]{Conjecture}
\newtheorem{definition}[theorem]{Definition}
\theoremstyle{remark}
\newtheorem{remark}[theorem]{Remark}

\newcommand{\R}{\mathbb R}
\newcommand{\wL}{L^{3,\infty}}

\title{The Averaging Barrier for the Radial-Scalar Q4 Budget Stack,\\
and the 2026-07-27 External Course Record}
\author{John Robert Lawson and Claude (Fable 5)}
\date{External audit round, 27 July 2026}

\begin{document}
\maketitle

\begin{abstract}
This round accompanies the 2026-07-27 external curation pass and states
the method-level finding behind its anti-drift rules.  Relative to
Tao's selected finite-time-blow-up averaged operator, an argument made
entirely from premises verified for both equations cannot prove global
regularity.  The charged \(q=4\) budgets are candidates for this
operator-relative classification, but the mixed Lorentz transfer
estimate still requires item-by-item verification.  Local energy, the
true quadratic pressure identity, physical-space carrier structure,
and exact triad geometry are natural places where the true equation may
break portability.  A shell realization would test only the finite
budget stack imposed on it; an averaged-equation realization would make
a stronger operator-relative statement.  Neither failure branch
automatically supplies a new portable obstruction.  The process record fixes two
discipline defects found in this audit: the four 2026-07-27 curation
commits were never pushed to either public remote, and their batch
omitted the required round report.  No Navier--Stokes theorem is proved
here, and the Clay problem is unsolved.
\end{abstract}

\section{Epistemic scope}

The imported blow-up theorem is a peer-reviewed result of
Tao~\cite{Tao16}.  The energy-measure theorem cited for contrast is the
peer-reviewed result of Leslie and Shvydkoy~\cite{LS18}.  The proposed
classification of this repository's budgets into operator-portable and
averaging-breaking is a same-system audit conjecture over the recorded
proofs; it awaits item-by-item transfer and external mathematical
review.  The observation that
regularity arguments insensitive to nonlinearity averaging cannot close
is Tao's own stated moral; no literature novelty is claimed for
Theorem~\ref{thm:barrier}.  The contribution of this round is the
file-anchored audit of the exact \(q=4\) stack, the resulting course
discipline, and the split realization obligations of
Section~\ref{sec:realization}.

Nothing here is a regularity theorem, a breakdown construction, an
energy-equality theorem, or a resolution of any Clay alternative.

\section{The averaged class and the barrier}

Tao~\cite{Tao16} constructs a bilinear operator \(\tilde B\), an
average of rotated, dilated, frequency-localised components of the
Euler bilinear operator \(B\), such that the system
\begin{equation}
 \partial_tu=\nu\Delta u+\tilde B(u,u),
 \qquad
 \nabla\cdot u=0
 \label{eq:averaged}
\end{equation}
admits Schwartz divergence-free data whose unique mild solution blows
up in finite time.  The operator \(\tilde B\) preserves the energy
cancellation \(\langle\tilde B(u,u),u\rangle=0\), hence the full energy
identity and energy class.  Its representation supplies standard
bilinear estimates with finite constants depending on the chosen
operator.  The endpoint Lorentz estimate used below is a transfer
obligation, not an imported theorem.

\begin{definition}[operator-portable premise]
\label{def:portable}
A premise is \emph{operator-portable relative to a fixed
\(\tilde B\)} if it has been proved for both the true Navier--Stokes
evolution and that \(\tilde B\) evolution.  Typical candidate inputs are:
(i) the energy identity and the energy-class norms it controls;
(ii) bilinear estimates on the nonlinearity proved with some finite
constant, which may depend on the selected operator; and
(iii) radial spectral bookkeeping, that is, calculus on pushforwards of
spectral densities under \(\xi\mapsto|\xi|\) and their heat transforms.
\end{definition}

\begin{theorem}[averaging barrier, immediate from~\cite{Tao16}]
\label{thm:barrier}
No valid global-regularity deduction can use only premises portable to
Tao's selected finite-time-blow-up operator.  At least one premise or
inference must distinguish the true Navier--Stokes nonlinearity from
that operator.
\end{theorem}

\begin{proof}
If every premise transferred, the same deduction would continue Tao's
blowing-up solution globally, a contradiction.
\end{proof}

\begin{remark}
Theorem~\ref{thm:barrier} neither classifies the current budgets nor
forbids closing one exact schedule such as the \(q=4\) record ledger.
Tao's known blow-up follows its own schedule, and the selected averaged
operator might avoid the \(q=4\) cell entirely.  Section~\ref{sec:realization}
therefore separates a shell-model test from the stronger
averaged-operator question.
\end{remark}

\section{Classification of the charged q4 stack}

The budgets actually charged against the \(q=4\) cell in the final
eight research rounds, with their recorded proofs, classify as follows.

\emph{Candidate operator-portable.}
(1) The energy identity and half-energy law \(A'=-AL\)
(rounds of 2026-07-25, radial triad spectrum and successors).
(2) \(M=\|u\|_{\wL}\in L^4_t\) by norm interpolation in the energy
class, and the schedule-conditional record clock \(M\in L^s_t\),
\(s<11/2\): pure arithmetic of the chosen ledger.
(3) \(Y=\|\nabla v\|_2\in L^2_t\): the energy inequality.
(4) The weighted transfer budget
\(\mathfrak q\lesssim\nu^{-1/2}\|v\|_2\|F\|_{\dot H^{-1}}\lesssim
\nu^{-1/2}\sqrt d\,MY\)
of the spectral quantile-speed round: Cauchy--Schwarz in frequency plus
Lorentz duality.  The corresponding mixed Lorentz estimate for the
selected averaged operator remains to be proved.
(5) The fixed-energy quantile speed law and its finite charge: kinematic
calculus on the radial density, provided the corresponding balance law
has transferred.
(6) The Gaussian, diagonal, early-reuse, and recycling ledgers: heat
transforms of radial spectra.

\emph{Candidate averaging-breaking inputs.}
(1) The Leslie--Shvydkoy energy-measure input~\cite{LS18} used by the
temporal five-power barrier and the terminal-dimension pincer: its
proof runs through the local energy balance of suitable solutions;
the rotations, dilations, and order-zero multipliers in \(\tilde B\)
do not preserve that local structure, so an analogue is not automatic
for~\eqref{eq:averaged}.
(2) The \(R^{5/2}\) carrier turnover clock: a physical-space statement
about energy in moving balls.
(3) The adjoint-pressure cost bounds of frozen ROUTE-R3B: after
projection an averaged equation may be represented with a scalar
pressure, but it need not retain the true local quadratic pressure
identity.
(4) Strong-trace and boundary-drift structure.

\begin{conjecture}
\label{prop:split}
After item-by-item transfer to Tao's selected operator, every budget
charged against the \(q=4\) cell through 2026-07-25 will be
operator-portable, while each milestone-registry result in rows 2--5
will retain at least one input not shared by that operator.
\end{conjecture}

The present evidence is narrower.  Explicit positive radial spectra
satisfy the finite scalar budget stack and show that those scalar
inequalities alone do not exclude the broad profile.  They are not
dynamics, do not settle future budgets, and do not prove that the
countermodel search was logically forced.

This is the formal content behind the router rules added on
2026-07-27: a new budget cycle must name structure that fails under
cascade averaging (local energy, pressure, or triad
phase/support geometry), and R3C's freeze counter restarted with the
radial-scalar budgets recorded as spent.

\section{Realization obligation}
\label{sec:realization}

\textbf{O-AVG-001a (open, shell rung).}  Decide whether a specified
shell-model class preserving the energy identity and an explicitly
enumerated finite budget stack admits a trajectory whose
radial spectral energy realises the broad-shell profile
\[
 e_\phi(\lambda,s)=\frac{A(s)}{L(s)}\,
 \phi\!\Bigl(\frac{\lambda}{L(s)}\Bigr),
 \qquad
 L(s)=s^{-9/11}\ell(s)^{2/11},
 \qquad
 A=\frac d2e^{-\int_0^sL},
\]
of the spectral quantile-speed round along the \(q=4\) schedule.  A
positive construction proves compatibility only with the budgets
actually verified on that model.  A negative result may be
model-specific and has no automatic route consequence.

\textbf{O-AVG-001b (open, averaged rung).}  Decide whether the selected
Tao averaged operator admits a trajectory in the same \(q=4\) cell.  A
positive construction gives an operator-relative barrier for any cell
exclusion.  A proof that this averaged evolution avoids the cell does
not automatically furnish a quantitative estimate portable back to the
true equation.

\begin{conjecture}
\label{conj:realizable}
The broad-shell profile is realisable in a specified shell-model class
preserving the energy identity and the present finite scalar budget
stack.
\end{conjecture}

Conjecture~\ref{conj:realizable} is not proved.  A realisation would
not be a Navier--Stokes solution and would resolve no Clay
alternative; it would bound the method only.

\section{Process record of this round}

\begin{enumerate}
\item \emph{Stranded commits.}  The four curation commits of
2026-07-27 (status scoping, milestone registry and router rules, the
audit record, the retirement) were found unpushed on both public
remotes, violating the repository's push-and-verify rule.  This round's
batch pushes \texttt{main} to \texttt{origin} and \texttt{brc} and
verifies the SHAs.
\item \emph{Missing round report.}  The curation batch omitted the
required self-contained round report.  This document restores the
discipline and records the omission; the underlying audit content is in
\texttt{dossier/reviews/2026-07-27-fable-audit.md}.
\item \emph{Stale provenance.}  An exhaustive reference sweep found
five prose citations of retired notes surviving inside three archival
rounds of 2026-07-24; the retirement criterion had resolved markdown
links only.  Each citation is now annotated in place with the Git
recovery recipe.  Imports, ledger artifact paths, markdown links, and
the Makefile were verified clean.
\item \emph{Freeze semantics.}  ROUTE-R3C's budget/countermodel cycle
count restarted at the 2026-07-27 curation.  Its next M4 cycle, after
Gate 0 verification, counts as cycle one of at most three.  Pure angular
mode counting is a portability candidate expected to admit countermodels; the exact
triad convolution support and phase constraints, local energy, and
pressure are the admissible inputs.
\item \emph{Verification queue.}  Before any new budget cycle, the
milestone registry rows are to be verified in order of external value:
the 2607 audit polished into an erratum-grade note; the
Leslie--Shvydkoy hypotheses of the five-power barrier checked against
the published text, including the Leray--Hopf continuation setting; the
\(R^{5/2}\) clock; the pincer together with the Albritton--Barker
amplitude reading.  External review remains the binding constraint:
same-system recomputation is exhausted as an error-finding instrument
for these results.
\end{enumerate}

\section{Open obligations}

\begin{itemize}
\item \textbf{O-AVG-001a/b}: test the broad-shell profile first in a
specified shell class and separately in Tao's selected averaged equation
(Section~\ref{sec:realization}).
\item \textbf{O-LS18-VERIFY}: verify the exact statement and scope
of~\cite{LS18} against the published text as used by milestone rows 2
and 4.
\item \textbf{O-2607-ERRATUM}: condense the 2607 audit into a
communicable erratum-grade comment note for external circulation.
\item \textbf{O-AB-AMPLITUDE}: confirm or retract the load-bearing
critical-amplitude reading of Albritton--Barker against the published
proof.
\item \textbf{O-ANGULAR-TRIAD}: the queued M4 cycle, restricted to
averaging-breaking inputs as recorded above.
\end{itemize}

\begin{thebibliography}{9}

\bibitem{Tao16}
T. Tao,
\emph{Finite time blowup for an averaged three-dimensional
Navier--Stokes equation},
Journal of the American Mathematical Society \textbf{29} (2016),
601--674.

\bibitem{LS18}
T.~M. Leslie and R. Shvydkoy,
\emph{The Energy Measure for the Euler and Navier--Stokes Equations},
Archive for Rational Mechanics and Analysis \textbf{230} (2018),
459--492.

\end{thebibliography}

\end{document}
